% =====================================================================
%  FORMULACION FINAL
%  Salida coronaria no-Newtoniana con bomba intramiocardica
%  Modelo acoplado LV 0D <-> coronaria 3D (CoupledLV0DCoronary3D)
%  Requiere: amsmath, amssymb, booktabs
% =====================================================================
\documentclass[11pt]{article}
\usepackage{amsmath,amssymb}
\usepackage{booktabs}
\usepackage[margin=2.3cm]{geometry}
\newcommand{\dd}[2]{\frac{\partial #1}{\partial #2}}
\title{Salida coronaria no-Newtoniana con bomba intramioc\'ardica:\\
       formulaci\'on final}
\date{}
\begin{document}\maketitle

\begin{abstract}
\noindent
Se presenta la formulaci\'on final de la condici\'on de salida coronaria del
modelo acoplado 0D--3D. Se conserva la ley de tubo WRMS evaluada con la misma
reolog\'ia de Carreau--Yasuda empleada en el dominio tridimensional, junto con
la compliancia expl\'icita y la bomba intramioc\'ardica. La revisi\'on afecta a
tres aspectos: la \emph{topolog\'ia}, con la compliancia situada en su posici\'on
microvascular; la \emph{interpretaci\'on anat\'omica} de los elementos; y la
\emph{estrategia de calibraci\'on}, basada en velocidades fisiol\'ogicas
prescritas. El n\'umero de constantes libres se reduce de 33 a 13, y todas las
restantes son magnitudes medibles.
\end{abstract}

\section{Principio rector}

\begin{quote}
\emph{La geometr\'ia representa anatom\'ia; la resistencia hidr\'aulica
representa la red vascular no resuelta.}
\end{quote}
En consecuencia: los radios medidos sobre la malla no se modifican; la
resistencia la produce el n\'umero de vasos en paralelo y la extensi\'on de la
red; y la reolog\'ia act\'ua a trav\'es de la ley constitutiva y no mediante
alteraciones de la geometr\'ia.

\section{Reolog\'ia y ley de tubo}

\begin{equation}
  \mu(\dot\gamma)=\mu_\infty+(\mu_0-\mu_\infty)
  \Bigl[1+(\lambda\dot\gamma)^{a}\Bigr]^{\frac{n-1}{a}} .
  \label{eq:cy}
\end{equation}
La condici\'on patol\'ogica se codifica mediante el conjunto completo
$(\mu_0,\mu_\infty,\lambda,n,a)$. La ecuaci\'on \eqref{eq:cy} es la misma en el
dominio 3D y en los elementos 0D, de modo que el acoplamiento es
constitutivamente consistente.

Para un tubo recto de radio $r$ y longitud $L$, con $\tau_w=r|\Delta p|/2L$ y
$\mu(\dot\gamma_w)\dot\gamma_w=\tau_w$,
\begin{equation}
  q(\Delta p;L,r)=\frac{\pi r^3}{\tau_w^{3}}
  \int_{0}^{\dot\gamma_w}\!\dot\gamma^{3}\,\mu^2\,
  \bigl(\mu+\dot\gamma\,\mu'\bigr)\,d\dot\gamma ,
  \qquad
  \boxed{\;\dd{q}{\Delta p}=\frac{\pi r^3\,\dot\gamma_w-3\,q}{\Delta p}\;}
  \label{eq:wrms}
\end{equation}
La derivada es exacta, por lo que el Jacobiano del m\'etodo de Newton es
consistente.

\paragraph{Relaci\'on de diagn\'ostico.} Con $Q=\pi r^2\bar v$,
\begin{equation}
  \dot\gamma_w=\frac{4Q}{\pi r^3}=\frac{4\bar v}{r},
  \label{eq:shear}
\end{equation}
independiente de $L$: el radio y la velocidad fijan el r\'egimen reol\'ogico,
mientras que la longitud s\'olo escala la resistencia. Esta identidad gobierna
el dise\~no de la calibraci\'on.

\section{Elementos y \'arboles en paralelo}

Cada componente es un elemento hidr\'aulico no lineal $Q=H(\Delta P;\Theta)$ con
$\Theta=(r,L,N,\mu,\dots)$. Un lecho de $N$ vasos id\'enticos en paralelo
satisface
\begin{equation}
  q_{\text{lecho}}(\Delta p)=N\,q(\Delta p;L,r),
  \qquad
  q_{\text{vaso}}=\frac{Q}{N} .
  \label{eq:paralelo}
\end{equation}
Esto permite conservar el calibre anat\'omico por vaso mientras la resistencia
la fija $N$. Un tubo \'unico no puede satisfacer ambas condiciones: al reducir
su radio para elevar $R$, la relaci\'on \eqref{eq:shear} conduce a velocidades
irreales y a $\mu\to\mu_\infty$.

\section{Topolog\'ia}
\label{sec:topo}

\[
\text{salida 3D}\;\rightarrow\;
\underbrace{\text{conducto}}_{r_p\ \text{de la malla}}\;\rightarrow\;
\underbrace{p_{ca},\;C_a}_{\text{epic\'ardico}}\;\rightarrow\;
\underbrace{\text{\'arbol arteriolar}}_{N_a}\;\rightarrow\;
\underbrace{p_c,\;C\;[\,p_{im}(t)\,]}_{\text{microvascular}}\;\rightarrow\;
\underbrace{\text{\'arbol venular}}_{N_v}\;\rightarrow\; P_{RA}
\]

\noindent
Corresponde a la disposici\'on cl\'asica de bomba intramioc\'ardica
$R_a-C(p_{im})-R_v$ con una compliancia arterial epic\'ardica antepuesta. La
compliancia microvascular $C$ se sit\'ua \emph{aguas abajo} de la resistencia
arteriolar (arteriolas terminales, capilares y v\'enulas), donde reside el
volumen sangu\'ineo coronario y sobre lo que act\'ua la compresi\'on
mioc\'ardica. La compliancia epic\'ardica $C_a$ se sit\'ua aguas arriba y, al
estar fuera del miocardio, no es comprimida por $p_{im}$.

\paragraph{Relevancia de la posici\'on.} La fuente $Q_{im}$ inyectada en $p_c$
se reparte entre las dos v\'ias en proporci\'on inversa a su resistencia. Si la
compliancia se sit\'ua aguas arriba del \'arbol arteriolar, las dos v\'ias son
el conducto epic\'ardico ($\sim\!10^{6}$) y el lecho arteriolar
($\sim\!10^{10}$~Pa\,s/m$^3$), de modo que pr\'acticamente todo el caudal de la
bomba se dirige hacia la arteria, produciendo picos retr\'ogrados de gran
amplitud y dejando inactivo el lado venoso. En la posici\'on microvascular las
v\'ias son el lecho arteriolar y el venular, cuyo cociente de resistencias es
del orden de $7$, lo que reparte el caudal de forma fisiol\'ogica entre un
impedimento sist\'olico arterial y una salida venosa sist\'olica.

\paragraph{Reparto de la ca\'ida de presi\'on.} La calibraci\'on debe
reproducir la distribuci\'on fisiol\'ogica de la p\'erdida de carga:

\begin{center}\small
\begin{tabular}{lcc}
\toprule
Segmento & Fisiol\'ogico & Modelo\\
\midrule
Conducto epic\'ardico & $\sim5\%$ & $<1\%$\\
Arteriolas y capilares & $\sim75\%$ & $87\%$\\
V\'enulas y venas & $10$--$15\%$ & $13\%$\\
\bottomrule
\end{tabular}
\end{center}

\noindent
Esto exige que el elemento venoso represente el \textbf{lecho venular} y no las
venas colectoras. La distinci\'on es decisiva: con calibre de vena colectora
($r_v\!\sim\!1$~mm) el compartimento venoso queda pr\'acticamente sin
resistencia ($0.26\%$ de la ca\'ida), lo que suprime simult\'aneamente el
impedimento sist\'olico, amortigua\-miento del pico venoso y el decaimiento
diast\'olico. Con calibre venular la constante de tiempo resultante,
$\tau=C(R_a\!\parallel\!R_v)$, pasa de $\sim\!4$~ms a $\sim\!0.2$~s, que es el
valor fisiol\'ogico del vaciado diast\'olico coronario.

\paragraph{Sobre la supresi\'on del segundo compartimento.} Conviene precisar
qu\'e se elimina y qu\'e se conserva. Se suprime el \emph{nodo de
almacenamiento} venoso ($C_v$, $p_{ven}$, $p_{im,v}$), no el \emph{elemento
hidr\'aulico} venoso. El \'arbol venoso, con su ley \eqref{eq:wrms} y su radio
ancho, se conserva \'ntegramente y pasa a ser la \'unica v\'ia de drenaje,
transportando el $99.7\%$ del caudal de la bomba frente a una fracci\'on
despreciable en la configuraci\'on anterior. Adem\'as, el nodo que subsiste es
precisamente la antigua uni\'on arteriolo-venosa: se elimin\'o el nodo situado
aguas arriba, que carec\'ia de correspondencia anat\'omica. La sensibilidad
reol\'ogica venosa, por tanto, no se reduce sino que aumenta.

\section{Presi\'on tisular y bomba intramioc\'ardica}

\begin{align}
  p_{im}^{\text{tgt}} &= \alpha\,p_{LV}+\alpha_a\,\tau_c,
  \label{eq:tgt}\\
  p_{im}^{n} &= p_{im}^{n-1}+\theta\bigl(p_{im}^{\text{tgt}}-p_{im}^{n-1}\bigr),
  \qquad \theta=\frac{\Delta t}{\Delta t+\tau_f},
  \label{eq:filt}\\
  Q_{im}&=C\,\frac{p_{im}^{n}-p_{im}^{n-1}}{\Delta t}.
  \label{eq:pump}
\end{align}
El retardo de primer orden \eqref{eq:filt} representa la velocidad finita de
desarrollo de la presi\'on tisular y acota $\dot p_{im}$ en las fases
isovolum\'etricas, donde $p_{LV}$ es casi discontinua.

\section{Balance de nodo y discretizaci\'on}

\paragraph{Almacenamiento colapsable.} El compartimento microvascular contiene
el volumen transmural $V(p_c-p_{im})$. Los vasos intramioc\'ardicos son de pared
delgada y se cierran cuando la presi\'on transmural se anula, de modo que $V$
debe ser \textbf{no negativo}: una compliancia lineal permite vaciar el lecho
m\'as all\'a de vac\'o, lo que hace caer $p_c$ por debajo de $P_{RA}$ y genera
reflujo venoso esp\'ureo. Se adopta
\begin{equation}
  V(p)=C\,p_0\,\ln\!\bigl(1+e^{p/p_0}\bigr),
  \qquad
  \frac{dV}{dp}(p)=\frac{C}{1+e^{-p/p_0}},
  \label{eq:vcol}
\end{equation}
que es lineal de pendiente $C$ donde el lecho est\'a distendido y satura a
volumen y compliancia nulos al colapsar. En el punto de operaci\'on
($p_c-p_{im}\approx1740$~Pa, varias veces $p_0$) se recupera
$dV/dp=0.997\,C$, de modo que la calibraci\'on no se altera.

\paragraph{Balance.} Escribir el balance en t\'erminos de $V$ y no de una
compliancia constante hace la bomba intramioc\'ardica \emph{impl\'cita} y
conserva el volumen exactamente:
\begin{equation}
  C_a\,\frac{dp_{ca}}{dt}=Q-q_a,
  \qquad
  \frac{d}{dt}V\bigl(p_c-p_{im}\bigr)=q_a-q_v .
  \label{eq:balance}
\end{equation}
Con Euler impl\'icito, inc\'ognitas $(p_{ca},p_c)$ y Newton $2\times2$,
\begin{align}
  R_1&=\frac{C_a}{\Delta t}\bigl(p_{ca}-p_{ca}^{n}\bigr)+q_a-Q=0,
  &q_a&=N_a\,q\bigl(p_{ca}-p_c;L_a,r_a\bigr),
  \label{eq:R1}\\
  R_2&=\frac{V(p_c-p_{im})-V^{n}}{\Delta t}-q_a+q_v=0,
  &q_v&=N_v\,q\bigl(p_c-P_{RA};L_v,r_v\bigr),
  \label{eq:R2}
\end{align}
\begin{equation}
  J=\begin{pmatrix}
    \dfrac{C_a}{\Delta t}+q_a' & -q_a'\\[6pt]
    -q_a' & \dfrac{1}{\Delta t}\dfrac{dV}{dp}+q_a'+q_v'
  \end{pmatrix},
  \qquad
  \boxed{\;p_{out}=p_{ca}+\Delta P_p(Q)+Z_c\,Q\;}
  \label{eq:jac}
\end{equation}
La formulaci\'on lineal se recupera haciendo $V=C\,p$, en cuyo caso
$R_2$ reproduce t\'ermino a t\'ermino la versi\'on con fuente expl\'cita
$Q_{im}=C\,\dot p_{im}$.

\paragraph{Impedancia caracter\'stica.} Un vaso compliante presenta
necesariamente una impedancia de onda, que en la salida vale
\begin{equation}
  Z_c=\frac{\rho\,c}{A},
  \qquad
  p_{out}=p_{ca}+\Delta P_p(Q)+Z_c\,Q ,
  \label{eq:zc}
\end{equation}
con $c$ la velocidad de onda de pulso y $A$ el \'area medida de la salida. Es el
$R_c$ del Windkessel de tres elementos. Su omisi\'on deja la compliancia
arterial y la inertancia del fluido del dominio 3D como un resonador
pr\'acticamente sin amortiguar, con periodo $2\pi\sqrt{L\,C_a}$ y factor de
calidad del orden de $30$, lo que produce oscilaciones espurias sostenidas en la
presi\'on de salida. No es un par\'ametro libre: se sigue de la densidad, del
\'area de la malla y de la velocidad de onda, y a lo largo de todo el rango
fisiol\'ogico ($c=5$--$10$~m/s) representa el $1.5$--$3\%$ de la resistencia de
rama y da un amortiguamiento $\zeta=0.26$--$0.51$.

\paragraph{Estabilidad del acoplamiento 3D--0D.} La presi\'on de salida se
aplica con un retraso de un paso, por lo que la resistencia que aparece en ella
queda tratada de forma expl\'cita. El esquema es estable si
\begin{equation}
  R_{\text{expl}}\;<\;Z_{3D}\approx\frac{\rho L}{A\,\Delta t},
  \label{eq:estab}
\end{equation}
condici\'on que la resistencia arteriolar, con el $87\%$ del presupuesto de
presi\'on, no satisface: exponerla directamente al dominio 3D hace divergir el
c\'alculo. La compliancia epic\'ardica resuelve el problema porque el 3D pasa a
ver un \emph{nodo capacitivo} en lugar de una resistencia dominante; entonces
\eqref{eq:estab} se sustituye por $\Delta t/C_a<Z_{3D}$, y ambas resistencias
quedan resueltas impl\'citamente en \eqref{eq:R1}--\eqref{eq:R2}. Con
$C_a=0.1\,C$ la condici\'on se satisface con m\'as de un orden de magnitud de
margen.

\section{Calibraci\'on}

Se ejecuta una sola vez con la viscosidad newtoniana de referencia $\mu_N$ y se
congela, de manera que un cambio de reolog\'ia modifica el flujo y nunca la
geometr\'ia, garantizando una comparaci\'on imparcial entre condiciones.

\paragraph{Paso 1: anatom\'ia.}
\begin{equation}
  r_{p,i}=\sqrt{A_i/\pi},
  \quad
  w_i=\frac{r_{p,i}^{3}}{\sum_j r_{p,j}^{3}},
  \quad
  Q_i=Q_{\text{tot}}\,w_i,
  \quad
  \Delta P_{p,i}=\frac{8\mu_N L_{p,i}Q_i}{\pi r_{p,i}^{4}} .
  \label{eq:paso1}
\end{equation}

\paragraph{Paso 2: \'arboles a partir de velocidades fisiol\'ogicas.}
\begin{equation}
  N_a=\frac{Q_i}{\bar v_a\,\pi r_a^{2}},
  \qquad
  N_v=\frac{Q_i}{\bar v_v\,\pi r_v^{2}},
  \qquad
  \dot\gamma_a=\frac{4\bar v_a}{r_a},
  \qquad
  \dot\gamma_v=\frac{4\bar v_v}{r_v}.
  \label{eq:paso2}
\end{equation}
La velocidad es un observable fisiol\'ogico directo y, por \eqref{eq:shear},
fija el punto de operaci\'on reol\'ogico. Los valores $N_v<1$ son admisibles y
representan la fracci\'on de una vena colectora compartida por varios
territorios.

\paragraph{Paso 3: presupuesto de presi\'on.} La vena posee radio y longitud
anat\'omicos, de modo que su ca\'ida es una salida del modelo; la extensi\'on
efectiva del \'arbol arteriolar absorbe el resto:
\begin{equation}
  \Delta P_v=\frac{8\mu_N L_v Q_v^{(1)}}{\pi r_v^{4}},
  \quad
  \Delta P_a=\Delta P-\Delta P_{p}-\Delta P_v,
  \quad
  L_a=\frac{\Delta P_a\,\pi r_a^{4}}{8\mu_N Q_a^{(1)}},
  \label{eq:paso3}
\end{equation}
con $Q_a^{(1)}=\bar v_a\pi r_a^2$, $Q_v^{(1)}=\bar v_v\pi r_v^2$ y
$\Delta P=p_{ar}(0)-\alpha p_{LV}(0)$. Todas las expresiones son cerradas, ya
que para $Q$ y $r$ dados $\tau_w$ queda determinado y $\Delta p\propto L$. La
longitud $L_a$ es un camino efectivo a trav\'es de varias generaciones en serie,
no un vaso individual.

\paragraph{Paso 4: compliancia medida.}
\begin{equation}
  C_i=C_{\text{tot}}\,w_i ,
  \qquad
  \tau=C\,\bigl(R_a\parallel R_v\bigr)\approx C\,R_v .
  \label{eq:paso4}
\end{equation}
La constante de decaimiento deja de ser un par\'ametro de ajuste y pasa a ser
una predicci\'on verificable.

\section{Constantes}

\begin{center}\small
\begin{tabular}{llll}
\toprule
S\'imbolo & C\'odigo & Valor & Origen\\
\midrule
$\mu_0,\mu_\infty$ & \texttt{viscosity.mu0,.muInf} & 0.353,\;4.181e-3 Pa\,s & reolog\'ia medida\\
$\lambda,n,a$ & \texttt{viscosity.lambda,.n,.yasuda} & 15.68 s,\;0.205,\;0.650 & reolog\'ia medida\\
$\mu_N$ & \texttt{newtonianCalibrationViscosity} & 3.5e-3 Pa\,s & referencia de calibraci\'on\\
\midrule
$r_p$ & \texttt{geometricParam.Rp} & de la malla & anatom\'ia (1.5--3.3 mm)\\
$L_p$ & \texttt{geometricParam.Lp} & 10--12.5 mm & anatom\'ia\\
$r_a$ & \texttt{arteriolarRadius} & 25 $\mu$m & anatom\'ia (arteriola)\\
$r_v$ & \texttt{venousRadius} & 30 $\mu$m & anatom\'ia (v\'enula)\\
$L_v$ & \texttt{venousLength} & 1.5 cm & camino venular efectivo\\
$\bar v_a$ & \texttt{arteriolarVelocity} & 5 mm/s & fisiolog\'ia (2--10 mm/s)\\
$\bar v_v$ & \texttt{venousVelocity} & 3 mm/s & fisiolog\'ia (1--5 mm/s)\\
$Q_{\text{tot}}$ & \texttt{lcaTargetFlow} & 2.0e-6 m$^3$/s & 120 mL/min\\
$C_{\text{tot}}$ & \texttt{coronaryComplianceTotal} & 3.2e-10 m$^3$/Pa & 0.043 mL/mmHg\\
$C_a/C$ & \texttt{arterialComplianceFraction} & 0.10 & compliancia epic\'ardica\\
$c$ & \texttt{pulseWaveVelocity} & 8 m/s & onda de pulso (5--10)\\
$p_0$ & \texttt{collapsePressureScale} & 300 Pa & colapso microvascular\\
$P_{RA}$ & \texttt{rightAtrialPressure} & 600 Pa & $\approx$4.5 mmHg\\
$\alpha,\alpha_a$ & \texttt{intramyocardial(Active)Fraction} & 0.65,\;0.01 & transmisi\'on tisular\\
$\tau_f$ & \texttt{intramyocardialFilterTau} & 0.02 s & desarrollo de $p_{im}$\\
\bottomrule
\end{tabular}
\end{center}

\noindent
Cantidades derivadas: $N_a\!\sim\!10^{4}$--$10^{5}$, $N_v\!\sim\!0.6$--$6.8$,
$L_a\!\approx\!46$~mm, $C_i$, $R_a$, $R_v$ y $\tau$. Total: 13 constantes
libres, todas ellas magnitudes medibles (radio, longitud, velocidad,
compliancia, caudal, presi\'on); ninguna es una resistencia equivalente.

\section{Sensibilidad reol\'ogica}

Con velocidades fisiol\'ogicas, el punto de operaci\'on arteriolar queda en
$\dot\gamma_a\!\approx\!800$~s$^{-1}$, donde la sangre es casi newtoniana. El
lecho venular opera a $\dot\gamma_v\!\approx\!400$~s$^{-1}$, pero su
cizallamiento recorre varias d\'ecadas a lo largo del ciclo:

\begin{center}\small
\begin{tabular}{lrr}
\toprule
Fase & $\dot\gamma_v$ (s$^{-1}$) & $\mu/\mu_\infty$\\
\midrule
S\'stole & 400 & 1.08\\
Di\'astole & 100 & 1.24\\
Nadir diast\'olico & 40 & 1.49\\
Bajo flujo & 10 & 2.43\\
Cuasi-estasis & 1 & 8.74\\
\bottomrule
\end{tabular}
\end{center}

\noindent
($\mu_0/\mu_\infty=84$ es el l\'mite de bajo corte.)

\noindent
La modulaci\'on reol\'ogica es, por tanto, considerable en t\'erminos
din\'amicos, aunque modesta si se eval\'ua \'unicamente en el punto medio. El
efecto de una condici\'on hiperviscosa debe buscarse en la forma de onda
durante las fases de bajo flujo (nadir diast\'olico, reversiones, salida venosa
sist\'olica) y no en el caudal medio, cuyo valor lo determina el lecho
arteriolar, que opera a alto cizallamiento y es pr\'acticamente newtoniano.

\section{Criterio para distinguir correcci\'on de ajuste}

Al a\~nadir un t\'ermino a un modelo cuyo comportamiento no satisface, conviene
someterlo a cuatro preguntas. Un t\'ermino leg\'timo las responde
afirmativamente; uno ajustado al resultado deseado, no.

\begin{enumerate}
\item \emph{¿Su valor procede de una medida independiente o de un primer
principio, o se fija a partir del resultado que se busca?}
\item \emph{¿El comportamiento es robusto en todo el rango fisiol\'ogico del
par\'ametro, o exige un valor de filo de navaja?}
\item \emph{¿Es defendible la ausencia del t\'ermino?} Si omitirlo es
f\'sicamente inconsistente, incluirlo es una correcci\'on, no un a\~nadido.
\item \emph{¿Se habr\'a incluido antes de observar el resultado?}
\end{enumerate}

\noindent
La impedancia caracter\'stica \eqref{eq:zc} las satisface: se deduce de
$\rho$, $A$ y $c$; el amortiguamiento es adecuado en todo el rango $c=5$--$10$
m/s; una compliancia sin impedancia de onda es f\'sicamente inconsistente; y
deber\'a haberse introducido junto con $C_a$.

La ley colapsable \eqref{eq:vcol} tambi\'en: lo que impone es la restricci\'on
$V\ge0$, no un valor; el punto de operaci\'on es insensible a $p_0$ en todo el
rango $150$--$600$~Pa, donde $dV/dp$ se mantiene entre $0.95\,C$ y $C$; extraer
de un compartimento m\'as sangre de la que contiene es imposible, luego su
ausencia es indefendible; y la restricci\'on deber\'a haberse impuesto desde el
planteamiento.

En cambio, elevar $\tau_f$ para reducir el pico de la bomba no las satisface:
carece de medida independiente que lo fije, se elegir\'a para alcanzar un
cociente objetivo y no corresponde al mecanismo real. Por ese motivo la
amplitud de la bomba se corrige imponiendo $V\ge0$ y no mediante ese
par\'ametro.

\section{Limitaciones}

\paragraph{Amplitud de la bomba.} Con la compliancia lineal, el volumen
desplazado por la s\'stole era $C\,\Delta p_{im}\approx3.2$~mL por latido, sin
cota superior, y la presi\'on transmural alcanzaba valores negativos
($\approx-4000$~Pa, esto es, $-1.3$~mL de volumen almacenado). La ley
colapsable \eqref{eq:vcol} acota la reserva expulsable al volumen realmente
contenido en reposo,
\begin{equation}
  V\bigl(p_c-p_{im}\bigr)\Big|_{\text{reposo}}
  =V(1740\ \text{Pa})\approx0.56\ \text{mL},
\end{equation}
es decir, un factor $\sim\!6$ menor, y garantiza $V\ge0$ en todo el ciclo. El
cociente pico/media del flujo venoso resultante es ahora una \emph{predicci\'on}
del modelo, y se reporta junto con las dem\'as m\'etricas de validaci\'on. Debe
observarse que la reserva expulsable queda fijada por $C$ y por la presi\'on
transmural de reposo, ambas determinadas de forma independiente, y no por un
par\'ametro ajustado a la forma de onda.

\paragraph{Presi\'on motriz.} Con la constante de tiempo del vaciado en su
rango fisiol\'ogico, la forma de onda diast\'olica depende tanto de $\tau$ como
de la evoluci\'on de la presi\'on motriz. Esta \'ultima la determinan
$p_{ar}(t)$, a trav\'es del Windkessel sist\'emico del modelo 0D, y $p_{im}(t)$,
a trav\'es de la onda de activaci\'on. Si el decaimiento observado siguiera
siendo demasiado r\'apido tras esta correcci\'on, dichos bloques son el lugar
donde debe continuar el an\'alisis.

\paragraph{Alcance del modelo reol\'ogico.} Los efectos no newtonianos de la
microcirculaci\'on asociados a la naturaleza particulada de la sangre
(F\aa hr\ae us--Lindqvist) no est\'an contenidos en \eqref{eq:cy} y requerir\'ian
un modelo espec\'ifico.

\end{document}
